%% 01-naskah-artikel-ilmiah.tex — bab dokumen contoh publikasi-laporan (hal. 57-58 (proposal), hal. 58-60 (laporan dan repositori)).
\chapter{NASKAH ARTIKEL ILMIAH}
Naskah artikel pada bagian ini disusun mengikuti template manuskrip jurnal tujuan, bukan
sistematika bab skripsi. Berkas terpisah naskah (berkas .docx atau .pdf) dilampirkan bersama
letter of acceptance, tautan homepage jurnal, dan bukti korespondensi dengan editor.

\section{Judul Artikel}
\section{Abstrak}
\section{Pendahuluan}
Naskah artikel memuat sitasi mengikuti gaya jurnal tujuan, misalnya:
Kajian pustaka disusun dari sumber terkini \citep{creswell2018research, patten2017understanding}
dan dari pedoman resmi universitas \citep{unnes2024panduanakademik}. Regulasi penjaminan mutu
menjadi dasar kebijakan \citep{permendikbudristek2023}.
 Contoh rujukan gaya MLA: \citemla[23]{patten2017understanding}.

\section{Metode}
\section{Hasil dan Pembahasan}
\section{Simpulan}
\section{Ucapan Terima Kasih}
